\documentclass[review,3p,times]{elsarticle}
\usepackage{amsmath,amsthm,verbatim,amssymb,amsfonts,dsfont,mathrsfs,amscd,graphicx}


%\usepackage[english]{babel}
\usepackage[russian, english]{babel} % Russian language support
\usepackage[T2A]{fontenc}   % Cyrillic font encoding
\usepackage{lmodern}        % Latin Modern font (supports Cyrillic)

\usepackage[utf8]{inputenc}
\usepackage{longtable}
\usepackage{pgfplots}
\usepackage{multirow}
\usepackage{makecell}
\usepackage{enumitem}
\usepackage{xfrac}
\usepackage[title]{appendix}

\usepackage{pdflscape}
\usepackage{hyperref}
\usepackage{lineno}
\usetikzlibrary{positioning}
\usepackage{listings}
\usepackage{xcolor}
\usepackage{graphicx}
\usepackage{float}

\usepackage{tikz}
\usetikzlibrary{shapes.geometric, arrows}
\usetikzlibrary{arrows.meta, positioning}

\lstset{basicstyle=\ttfamily, mathescape=true, escapechar=\%}

%\linenumbers

\usepackage[nodisplayskipstretch]{setspace}

\setenumerate{itemsep=0.02cm, topsep=0.05cm}
\setlist[itemize]{itemsep=0.02cm, topsep=0.05cm}

% Bibliography
\bibliographystyle{elsarticle-num}

\usepackage[tableposition=top]{caption}  % caption is on top of table
\usepackage{subcaption}
\captionsetup{skip=5pt}  % distance between tabular caption and tabular


\usepackage{tikz}
\usetikzlibrary{arrows,shapes,positioning,shadows,trees,calc}

\tikzset{
    basic/.style  = {draw, text width=2cm, drop shadow, font=\sffamily, rectangle},
    root/.style   = {basic, rounded corners=2pt, thin, align=center,
                     fill=green!30},
    level 2/.style = {basic, rounded corners=6pt, thin,align=center, fill=green!60,
                     text width=8em},
    level 3/.style = {basic, thin, align=center, fill=pink!60, text width=9.5em}
}


\usepackage[linesnumbered,ruled]{algorithm2e}
% Smaller font for algorithm comments
\newcommand\mycommfont[1]{\footnotesize\ttfamily\textcolor{black}{#1}}
\SetCommentSty{mycommfont}

% Basic math notation
\def\R{\mathbb R}
\def\E{\mathbb E}
\def\EE{\mathcal E}
\def\N{\mathbb N}
\def\NN{\mathcal N}
\def\NNN{\mathscr N}
\def\D{\mathbb D}
\def\DD{\mathcal D}
\def\DDD{\mathscr D}
\def\X{\mathbb X}
\def\XX{\mathcal X}
\def\XXX{\mathscr X}
\def\S{\mathbb S}
\def\SS{\mathcal S}
\def\T{\mathcal T}
\def\A{\mathcal A}
\def\H{\mathcal H}
\def\I{\mathcal I}
\def\O{\mathcal O}
\def\U{\mathcal U}
\def\F{\mathbb F}
\def\FF{\mathcal F}
\def\FFF{\mathscr F}
\def\L{\mathbb L}
\def\LL{\mathcal L}
\def\LLL{\mathscr L}

\newcommand{\Rho}{\mathrm{P}}

% A nice-looking symbol for the empty set
\let\oldemptyset\emptyset
\let\emptyset\varnothing

% A nice-looking argmin(-max) function
\DeclareMathOperator*{\argmax}{arg\,max}
\DeclareMathOperator*{\argmin}{arg\,min}

% A nice-looking epsilon symbol
\newcommand{\eps}{\varepsilon}

% Keywords command
\providecommand{\keywords}[1]
{
  \small	
  \textbf{\textit{Keywords:}} #1
}

\usepackage{titlesec}
% Section/subsection styles for experiments
% Syntax: \titlespacing{command}{left spacing}{before spacing}{after spacing}[right]
\newcommand\appxsection{%
  \titlespacing\section{0cm}{0cm plus 0cm minus 0cm}{0cm plus 0cm minus 0cm}%
  \titleformat*{\section}{\bfseries\fontsize{12}{14}\selectfont}
}
\newcommand\appxsubsection{%
  \titlespacing\subsection{0cm}{-0.3cm plus 0cm minus 0cm}{0cm plus 0cm minus 0cm}%
  \titleformat*{\subsection}{\bfseries\fontsize{10}{12}\selectfont}
}
% Smaller captions for experiments (footnotesize=10pt)
\newcommand\appxcaptions{%
  \captionsetup{font=footnotesize}%
  \captionsetup{skip=1pt}
}
\newlength{\tableheight}  % height of each table in experiments
\setlength{\tableheight}{1.2cm}
\newlength{\negspace}
\setlength{\negspace}{-0.25cm}  % space to bring tabulars up a bit
% Font size for experiments
\newcommand\appxfont{%
  \fontsize{10}{12}%
  \selectfont
}
\setlength\medskipamount{0.15cm plus 0cm minus 0cm}  % space between tabulars in experiments

\newtheorem{theorem}{Theorem}[section]
\newtheorem{lemma}[theorem]{Lemma}
\newtheorem{definition}{Definition}[section]

\pgfplotsset{compat=1.17}

\journal{Pattern Recognition}

\begin{document}

\begin{frontmatter}

\title{Информационно-разведочный поиск на основе метода количественной оценки концептуальной новизны научных публикаций}



\begin{keyword}
Exploratory Information Search \sep Novelty Assessment \sep Large Language Model \sep Clustering \sep Embedding \sep Knowledge Field \sep Breakthrough Potential Index (BPI) \sep Emerging Topics

\end{keyword}





\end{frontmatter}

\section{Аннотация}


В статье предлагается метод информационно-разведочного поиска научно-технических публикаций на основе количественной оценки их концептуальной новизны. Публикации рассматриваются как конфигурации ограниченного множества базовых идей, формирующих идейное пространство научно-технической информации. Для извлечения и структурирования этих идей используются семантическая сегментация текстов, эмбеддинговые представления и большие языковые модели. На основе идейного спектра публикации вводятся показатели новизны, редкости и когерентности её концептуальной конфигурации. Предлагаемый подход может применяться для выявления публикаций с высоким потенциалом научного прорыва и интеллектуального ранжирования научно-технической информации.



\section{Introduction}

Современный рост объёмов открытой научно-технической информации существенно осложняет задачи поиска, отбора и аналитической интерпретации публикаций, представляющих наибольший научный и технологический интерес. Традиционные методы информационного поиска в основном ориентированы на тематическое соответствие документов запросу, однако они в меньшей степени приспособлены к выявлению публикаций, содержащих концептуально новые идеи и обладающих потенциалом научного прорыва. В условиях высокой плотности информационных потоков особую значимость приобретает разработка методов, позволяющих оценивать не только релевантность публикаций, но и степень их концептуальной новизны.

В данной работе предлагается подход к информационно-разведочному поиску научно-технических публикаций, основанный на представлении информационного пространства как идейного поля, а отдельных публикаций — как конфигураций ограниченного множества базовых идей. Для построения такого представления используются методы семантической сегментации текста, эмбеддингового анализа и большие языковые модели, позволяющие извлекать и обобщать концептуальное содержание публикаций. На этой основе вводятся количественные показатели, характеризующие структурную новизну, редкость и когерентность идейной конфигурации публикации.

Цель исследования состоит в разработке метода количественной оценки концептуальной новизны научных публикаций и его использовании в задачах информационно-разведочного поиска. Предлагаемый подход ориентирован на выявление публикаций, находящихся в нетипичных, но концептуально согласованных областях идейного пространства, и тем самым может служить основой для интеллектуальных систем ранжирования, мониторинга и аналитического сопровождения научно-технической информации.



\section{Идейное пространство открытой научно-технической информации как объект исследования}

Открытая научно-техническая информация представляет собой не просто совокупность документов, сообщений и наблюдаемых информационных потоков, а более сложную среду, в которой циркулируют, взаимодействуют и преобразуются идеи. В этой среде научные публикации и патенты выполняют функцию не только текстовых носителей сведений, но и институционализированных форм фиксации, передачи и воспроизводства научно-технических концепций, интерпретационных рамок и технологических решений. Поэтому анализ информационного пространства целесообразно вести не только на уровне документов и их атрибутов, но и на уровне тех идейных структур, которые через эти документы проявляются.

Такой подход позволяет рассматривать информационное пространство как среду, в которой идеи распространяются, усиливаются, ослабевают, вступают в конкуренцию, сочетаются друг с другом и переходят в новые формы выражения. Его динамика определяется не только возникновением принципиально новых идей, но и перераспределением внимания между уже существующими концепциями, изменением их значимости, рекомбинацией и переинтерпретацией в новых научных и технологических контекстах. Иначе говоря, эволюция научно-технической информации во многом является отражением эволюции конфигураций идей, лежащих в её основе.

В этой связи информационное пространство общества может быть интерпретировано как \emph{поле знаний} (\emph{Knowledge Field}) или, в более узком смысле, как \emph{идейное поле} (\emph{Ideational Field}). Под идейным полем далее будем понимать распределённую во времени и пространстве систему концептов, интерпретационных схем, смысловых напряжений и их взаимных влияний, проявляющихся через тексты, события, коммуникации и социальные практики. В такой постановке отдельный текст выступает не как изолированный объект, а как локальная наблюдаемая проекция более общей структуры распределения и взаимодействия идей.

В рамках настоящего исследования в качестве наблюдаемой части такого идейного поля рассматриваются тексты открытых научно-технических публикаций. Это позволяет принять следующее базовое допущение: значительная часть публикаций представляет собой различные варианты выражения сравнительно ограниченного множества устойчивых и концептуально значимых научно-технических идей. Различия между текстами в этом случае обусловлены прежде всего не самим существованием полностью уникальных смыслов в каждом документе, а различиями в способах представления идей, в их сочетаниях, акцентах и интенсивности выражения. При таком взгляде множество конкретных публикаций оказывается значительно более изменчивым, чем лежащее в их основе множество базовых идей.

Исходя из этого формулируется основная гипотеза исследования: научные публикации могут быть представлены как конфигурации ограниченного множества базовых идей, а концептуальная новизна публикации определяется редкостью и когерентностью её конфигурации в пространстве этих идей.

Для формализации данного подхода введём множество базовых идей
\[
\mathcal{I}=\{I_1,\dots,I_K\},
\]
где каждая идея $I_i$ соответствует концепции, утверждению, интерпретационной рамке или устойчивому смысловому конструкту, обладающему научно-технической значимостью. Тогда каждой публикации $w$ может быть сопоставлен её идейный спектр
\[
\mathbf{s}(w)=(s_1,\dots,s_K),
\]
где $s_i$ характеризует степень выраженности идеи $I_i$ в данной публикации. На практике вектор $\mathbf{s}(w)$ является разреженным, поскольку отдельный документ, как правило, актуализирует лишь небольшую часть полного множества возможных идей.

Такое представление позволяет трактовать публикацию как локальную спектральную реализацию идейного поля. Если поле знаний понимается как глобальная структура распределения и взаимодействия идей, то идейный спектр отдельного текста можно рассматривать как наблюдаемую проекцию этой структуры в конкретной точке информационного пространства и в конкретный момент времени. Тем самым возникает следующая концептуальная последовательность:
\[
\text{идейное поле}
\;\rightarrow\;
\text{идейный спектр текста}
\;\rightarrow\;
\text{структурированное представление знаний}.
\]

Ключевой методологической задачей при таком подходе становится извлечение идейного спектра текста, то есть переход от неструктурированного текстового описания к формализованному представлению его концептуального содержания. В качестве интеллектуального оператора такого анализа могут использоваться большие языковые модели (LLM), обладающие способностью выполнять семантическую интерпретацию текста, выявлять его концептуальные элементы и соотносить локальные высказывания с более общими смысловыми структурами.

Формально операцию извлечения идейного спектра можно представить в виде
\begin{equation}
\mathbf{s}(w)=f_{\mathrm{LLM}}(w,p),
\label{eq:llm_operator}
\end{equation}
где $f_{\mathrm{LLM}}$ обозначает оператор интеллектуального концептуального анализа и интерпретации текста, а параметр $p$ задаёт постановку задачи извлечения идейно-концептуальных элементов и контекст её выполнения (prompt в терминах prompt-инженеринга).

Таким образом, в предлагаемой модели тексты открытых научно-технических публикаций рассматриваются как наблюдаемые проявления динамического идейного поля, а их концептуальный анализ — как средство реконструкции структуры этого поля. Такое представление создаёт основу для дальнейшего количественного описания идейных конфигураций публикаций и оценки их концептуальной новизны.



\section{Формирование множества базовых идей}

Множество базовых идей 
\[
\mathcal{I}=\{I_1,\dots,I_K\}
\]
может формироваться автоматически на основе корпуса текстов. Пусть 
\[
\mathcal{W}=\{w_1,\dots,w_N\}
\]
обозначает корпус документов.

\subsection{Сегментация текстов}

На первом этапе каждый документ $w \in \mathcal{W}$ сегментируется на последовательность фрагментов, обладающих относительным смысловым единством. Формально документ представляется как последовательность предложений
\[
w=(s_1,s_2,\dots,s_T).
\]

Для каждого предложения $s_t$ вычисляется его семантическое векторное представление (эмбеддинг)
\[
\mathbf{e}_t = E(s_t),
\]
где $E(\cdot)$ обозначает оператор получения эмбеддинга предложения с использованием большой языковой модели или специализированной модели семантических представлений.

Полученная последовательность эмбеддингов
\[
(\mathbf{e}_1,\mathbf{e}_2,\dots,\mathbf{e}_T)
\]
может рассматриваться как многомерный временной ряд в семантическом пространстве.

\subsection{Определение границ смысловых сегментов}

Границы смысловых сегментов определяются методом сегментации временных рядов по локальному максимуму разрыва (segmentation of time series by local maximum gap). Для каждой пары соседних предложений вычисляется расстояние

\[
d_t = \|\mathbf{e}_{t+1}-\mathbf{e}_t\|,
\]

где $\|\cdot\|$ обозначает выбранную метрику в пространстве эмбеддингов (например, косинусное расстояние).

Пусть $\{d_t\}_{t=1}^{T-1}$ — последовательность межпредложенческих расстояний. Для этой последовательности вычисляются статистические характеристики

\[
\mu_d = \frac{1}{T-1}\sum_{t=1}^{T-1} d_t, 
\qquad
\sigma_d = \sqrt{\frac{1}{T-1}\sum_{t=1}^{T-1}(d_t-\mu_d)^2}.
\]

Адаптивный порог семантического разрыва определяется как

\[
\tau = \mu_d + \alpha \sigma_d,
\]

где параметр $\alpha$ задаёт чувствительность процедуры сегментации.

Позиция $t$ рассматривается как граница сегмента, если выполняются условия локального максимума и порогового превышения:

\[
d_t > d_{t-1}, \qquad 
d_t > d_{t+1}, \qquad 
d_t > \tau .
\]

Иными словами, границами сегментов считаются только такие локальные максимумы последовательности $\{d_t\}$, величина которых превышает адаптивный порог $\tau$. Такие точки интерпретируются как позиции существенного семантического разрыва между соседними предложениями.

В результате документ представляется в виде последовательности смысловых сегментов

\[
w=(c_1,c_2,\dots,c_M),
\]

где каждый сегмент $c_j$ объединяет последовательность предложений, характеризующихся относительной семантической однородностью.



\subsection{Извлечение множества базовых идей}

На предыдущем этапе каждый документ корпуса $\mathcal{W}=\{w_1,\dots,w_N\}$ был представлен последовательностью смысловых сегментов $w=(c_1,c_2,\dots,c_M)$, где $c_j$ обозначает семантически однородный текстовый сегмент.



\subsubsection{Извлечение идейных формулировок сегментов}

Для каждого сегмента $c_j$ с использованием интеллектуального оператора концептуального анализа текста (\ref{eq:llm_operator}) выполняется извлечение его идейного содержания. При этом при анализе желательно учитывать не только сам сегмент $c_j$, но и его контекст. В идеальном случае таким контекстом выступает весь документ $w$, в котором данный сегмент содержится, поскольку интерпретация локального текстового фрагмента часто существенно зависит от общей логики и концептуальной структуры работы.

С учётом контекста операция извлечения идеи может быть представлена как $i_j = f_{\mathrm{LLM}}(c_j, w, p)$, где $i_j$ представляет собой текстовую формулировку основной идеи сегмента, $w$ обозначает контекст документа, а параметр $p$ задаёт соответствующее задание (prompt) для языковой модели.

Задание $p$ формулируется таким образом, чтобы извлекались преимущественно высокоуровневые и абстрактные формулировки идей. Основной акцент при этом делается на выявление концептуальной сути высказывания при минимизации конкретных деталей, частных примеров и технических уточнений.

Пример возможной формулировки задания $p$:

\begin{quote}
``Сформулируйте основную концептуальную идею данного текстового фрагмента с учётом контекста всего документа в виде одного краткого и максимально обобщённого утверждения. Сосредоточьтесь на концептуальной сути идеи, опуская частные детали, примеры и технические реализации.''
\end{quote}

В результате каждому текстовому сегменту сопоставляется компактная идейная формулировка, отражающая его основное концептуальное содержание.



\subsubsection{Эмбеддинговое представление идей}

Для каждой полученной идейной формулировки $i_j$ далее вычисляется её векторное семантическое представление

\[
\mathbf{v}_j = E(i_j),
\]

где $E(\cdot)$ обозначает оператор получения эмбеддинга текста. Полученный вектор $\mathbf{v}_j$ используется для анализа семантической близости идей.

\subsubsection{Формирование множества базовых идей}

Формирование множества базовых идей 
\[
\mathcal{I}=\{I_1,\dots,I_K\}
\]
осуществляется инкрементально по мере обработки всех сегментов корпуса.

Каждая базовая идея $I_k$ рассматривается как представитель кластера семантически близких идейных формулировок. Пусть $\mathbf{u}_k$ обозначает эмбеддинг представительной формулировки идеи $I_k$.

Для каждой новой извлечённой идеи $i_j$ с эмбеддингом $\mathbf{v}_j$ вычисляется расстояние до всех уже сформированных базовых идей

\[
d_k = \|\mathbf{v}_j - \mathbf{u}_k\|, 
\qquad k=1,\dots,K.
\]

После этого определяется минимальное расстояние

\[
d_{\min} = \min_{k} d_k,
\]

и соответствующая ему базовая идея

\[
k^* = \arg\min_{k} d_k.
\]

Далее применяется пороговое правило принятия решения. Пусть $\delta$ обозначает заданный порог семантической близости.

\begin{itemize}

\item Если
\[
d_{\min} < \delta,
\]
то новая идея $i_j$ считается семантически близкой к уже существующей базовой идее $I_{k^*}$ и добавляется в соответствующий кластер.

\item Если
\[
d_{\min} \ge \delta,
\]
то идея $i_j$ интерпретируется как концептуально новая и добавляется в множество базовых идей в качестве нового элемента

\[
I_{K+1} = i_j,
\qquad
K \leftarrow K+1.
\]

\end{itemize}

\subsubsection{Обновление представительного описания кластера}

Если новая идея добавляется в уже существующий кластер $I_{k^*}$, выполняется обновление его представительного описания.

Пусть
\[
\mathcal{C}_{k^*} = \{i^{(1)}, i^{(2)}, \dots, i^{(m)}\}
\]
обозначает множество всех идейных формулировок, входящих в данный кластер после добавления новой идеи.

Для этого множества выполняется обобщение (суммаризация) с использованием интеллектуального оператора

\[
I_{k^*} = f_{\mathrm{LLM}}(\mathcal{C}_{k^*}, p_{\mathrm{sum}}),
\]

где $p_{\mathrm{sum}}$ задаёт задачу обобщения множества близких идейных формулировок в единую более абстрактную концептуальную формулировку.

Пример возможной формулировки задания $p_{\mathrm{sum}}$:

\begin{quote}
``Обобщите приведённые идейные формулировки в одно краткое и максимально абстрактное утверждение, отражающее их общую концептуальную суть. Исключите частные детали и сохраните только наиболее общий смысл, объединяющий все формулировки.''
\end{quote}


Полученное обновлённое описание базовой идеи затем повторно преобразуется в эмбеддинг

\[
\mathbf{u}_{k^*} = E(I_{k^*}),
\]

который используется при дальнейшем сравнении идей.

\subsubsection{Интерпретация результата}

В результате обработки всех сегментов корпуса формируется множество базовых идей

\[
\mathcal{I}=\{I_1,\dots,I_K\},
\]

каждая из которых представляет собой обобщённое концептуальное описание кластера близких идейных формулировок, извлечённых из текстов корпуса.

Таким образом, процедура автоматически выявляет устойчивые концептуальные структуры, присутствующие в информационном пространстве научно-технических публикаций. Полученное множество базовых идей может рассматриваться как приближённое дискретное представление структуры идейного поля, формируемого данным корпусом текстов.




\subsection{Идейное представление документов в пространстве базовых идей}

В результате предыдущих этапов для каждого документа корпуса выполняется сегментация текста и извлечение идейных формулировок соответствующих сегментов. Таким образом, каждому документу $w$ сопоставляется множество идейных формулировок, соответствующих множеству составляющих его текстовых сегментов.

Поскольку на предыдущем этапе уже сформировано множество базовых идей
$\mathcal{I}=\{I_1,\dots,I_K\}$, можно построить для каждого документа его представление в пространстве этих идей. Для этого каждому документу $w$ сопоставляется вектор размерности $K$

\[
\mathbf{s}(w)=(s_1(w),\dots,s_K(w)),
\]

где компонент $s_i(w)$ количественно характеризует степень выраженности базовой идеи $I_i$ в документе $w$.

Для вычисления этих величин используются эмбеддинговые представления идейных формулировок сегментов документа и базовых идей. Пусть $\mathcal{J}(w)=\{i_1,\dots,i_M\}$ обозначает множество идейных формулировок, извлечённых из сегментов документа $w$, а $\mathbf{v}_j = E(i_j)$ — их эмбеддинги. Аналогично, пусть $\mathbf{u}_i = E(I_i)$ обозначает эмбеддинг базовой идеи $I_i$.

Тогда степень выраженности идеи $I_i$ в документе может быть оценена на основе семантической близости между эмбеддингом базовой идеи и эмбеддингами идей сегментов документа. В частности, одним из возможных определений является

\[
s_i(w)=\max_{i_j \in \mathcal{J}(w)} \operatorname{sim}(\mathbf{u}_i,\mathbf{v}_j),
\]

где $\operatorname{sim}(\cdot,\cdot)$ обозначает выбранную меру семантической близости (например, косинусное сходство).

Таким образом, для каждого документа корпуса формируется вектор его идейного представления в пространстве базовых идей. Иными словами, каждый текст корпуса отображается в пространство базовых идей, сформированное на основе всего корпуса документов. Такое представление позволяет рассматривать тексты как точки в $K$-мерном идейном пространстве и далее применять к ним методы количественного анализа, сравнения и кластеризации.





\section{Количественная оценка концептуальной новизны}

После построения пространства базовых идей каждая публикация $w$ может быть представлена в виде точки в $K$-мерном пространстве идейных спектров

\[
\mathbf{s}(w)\in \mathbb{R}^K.
\]

Соответственно, корпус публикаций

\[
\mathcal{W}=\{w_1,\dots,w_N\}
\]

образует облако точек в пространстве базовых идей. В этой интерпретации обычные публикации располагаются в областях высокой плотности идейного пространства, отражающих типичные конфигурации научных идей. Публикации, содержащие новые концептуальные решения, отклоняются от таких областей.

Концептуальная новизна публикации характеризует степень её нетипичности относительно существующего корпуса. Однако высокая новизна сама по себе ещё не означает научный прорыв: она может возникать и у маргинальных, случайных или шумовых работ. Поэтому при оценке новизны необходимо учитывать не только редкость идей, но и степень осмысленности и структурной когерентности их комбинации. В рамках предложенной модели научный прорыв интерпретируется как редкая, но концептуально согласованная конфигурация идей.

Для количественной оценки этих свойств могут использоваться несколько взаимодополняющих метрик.

\subsection{Структурная новизна}

Пусть

\[
\bar{\mathbf{s}}=\frac{1}{N}\sum_{n=1}^{N}\mathbf{s}(w_n)
\]

обозначает центр идейного поля, а

\[
\Sigma=\mathrm{Cov}(\mathbf{s}(w))
\]

— ковариационную матрицу идейных спектров корпуса.

Тогда структурная новизна публикации может быть определена как расстояние Махаланобиса

\[
D_M(w)=
\sqrt{
(\mathbf{s}(w)-\bar{\mathbf{s}})^T
\Sigma^{-1}
(\mathbf{s}(w)-\bar{\mathbf{s}})
}.
\]

Данная величина характеризует степень отклонения публикации от типичной структуры идейного поля с учётом корреляций между идеями. В отличие от обычного евклидова расстояния, эта метрика отражает именно необычность комбинации идей относительно общей структуры корпуса.

\subsection{Редкость идейной конфигурации}

Пусть

\[
p_i=\frac{1}{N}\sum_{n=1}^{N}s_i(w_n)
\]

обозначает среднюю выраженность идеи $I_i$ в корпусе публикаций.

Тогда информационная редкость идейной конфигурации публикации может быть оценена величиной

\[
R(w)=-\sum_{i=1}^{K}s_i(w)\log p_i.
\]

Эта метрика возрастает, если публикация опирается на идеи, которые редко встречаются в корпусе.

Однако прорыв обычно связан не просто с редкостью идей, но и с их концептуальной согласованностью. Для оценки структурной когерентности идейного профиля введём меру концентрации спектра

\[
C(w)=\sum_{i=1}^{K}s_i(w)^2.
\]

Если спектр $\mathbf{s}(w)$ нормирован, высокая величина $C(w)$ означает, что публикация формирует достаточно чёткий и структурированный идейный профиль, а не распределяется равномерно по множеству идей.

\subsection{Редкость положения в идейном пространстве}

Дополнительную характеристику новизны можно получить на основе анализа локальной плотности идейного пространства. Пусть $\mathcal{N}_k(w)$ обозначает множество $k$ ближайших публикаций к $w$ в пространстве идейных спектров. Тогда локальная плотность может быть оценена как

\[
\rho(w)=
\left(
\frac{1}{k}\sum_{w'\in \mathcal{N}_k(w)}
\|\mathbf{s}(w)-\mathbf{s}(w')\|
\right)^{-1}.
\]

Чем меньше значение $\rho(w)$, тем менее населена область идейного пространства, в которой расположена публикация.

\subsection{Индекс потенциала научного прорыва}

Используя перечисленные характеристики, можно определить интегральную оценку потенциала научного прорыва публикации — \emph{индекс потенциала научного прорыва}. Данный индекс отражает одновременно редкость идей, когерентность их комбинации и нетипичность положения публикации в идейном пространстве.

В простейшем случае индекс может быть определён как

\[
B_{\mathrm{pot}}(w)=D_M(w)\cdot C(w),
\]

где учитывается сочетание структурной новизны и когерентности идейного профиля.

Более общая форма индекса может учитывать также редкость идей и локальную плотность идейного пространства:

\[
B_{\mathrm{pot}}(w)=R(w)\cdot C(w)\cdot \frac{1}{\rho(w)}.
\]

В данной интерпретации

\begin{itemize}
\item $R(w)$ характеризует редкость используемых идей,
\item $C(w)$ отражает когерентность идейного профиля публикации,
\item $1/\rho(w)$ указывает на расположение публикации в малонаселённой области идейного пространства.
\end{itemize}

Таким образом, индекс $B_{\mathrm{pot}}(w)$ может рассматриваться как количественная оценка потенциала научного прорыва — способности публикации формировать новую концептуальную конфигурацию идей в структуре научного знания.






\section{Экспериментальная постановка}

Для экспериментальной проверки эффективности предложенного подхода целесообразно использовать две взаимодополняющие постановки. Первая ориентирована на содержательную валидацию метода и позволяет проверить, способен ли он выделять публикации, которые задним числом рассматриваются как концептуально значимые или прорывные. Вторая направлена на анализ вклада отдельных компонентов метода и позволяет установить, какие именно элементы предложенной модели обеспечивают итоговое качество оценки концептуальной новизны.

\subsection{Ретроспективное выявление известных прорывных публикаций}

Цель данной постановки состоит в проверке гипотезы о том, что публикации, признанные впоследствии значимыми для развития соответствующего научно-технического направления, должны получать высокие значения показателя концептуальной новизны, вычисленного на основе только предшествующего им корпуса публикаций.

Для проведения эксперимента формируется предметно однородный корпус научных публикаций за некоторый временной интервал
\[
\mathcal{W}=\{w_1,\dots,w_N\},
\]
упорядоченный по времени публикации. Внутри корпуса выделяется подмножество
\[
\mathcal{B}\subset \mathcal{W},
\]
содержащее публикации, которые на основе внешних критериев рассматриваются как известные прорывные работы. В качестве таких критериев могут использоваться экспертная разметка, включение публикации в обзорные работы, высокая отложенная цитируемость, широкое признание в профессиональном сообществе либо принадлежность к числу ранних работ, инициировавших быстрое развитие нового направления.

Для каждой публикации $w_t$ показатель концептуальной новизны вычисляется в ретроспективной постановке, то есть только относительно множества публикаций, предшествующих ей по времени:
\[
\mathcal{W}_{<t}=\{w_j \in \mathcal{W}\mid \mathrm{time}(w_j)<\mathrm{time}(w_t)\}.
\]
Это условие принципиально важно, поскольку исключает использование будущей информации при оценке новизны и тем самым обеспечивает корректность проверки. На основе корпуса $\mathcal{W}_{<t}$ формируется множество базовых идей, строится идейное пространство, после чего для публикации $w_t$ вычисляется её идейный спектр и итоговый индекс концептуальной новизны или потенциала научного прорыва:
\[
B_{\mathrm{pot}}(w_t).
\]

После вычисления значений $B_{\mathrm{pot}}(w)$ для всех публикаций анализируется положение элементов множества $\mathcal{B}$ в общем ранжировании корпуса. Эффективность метода может оцениваться по нескольким показателям: среднему рангу прорывных публикаций, доле таких публикаций в верхних $k$ позициях ранжирования, а также по стандартным метрикам ранжирования, например Precision@$k$ или nDCG@$k$, если для корпуса доступна соответствующая бинарная или градуированная разметка.

Ожидаемый результат эксперимента состоит в том, что публикации из множества $\mathcal{B}$ будут статистически значимо концентрироваться в верхней части ранжированного списка. Это будет свидетельствовать о том, что предложенный подход действительно способен выявлять публикации с нетипичными, но концептуально согласованными идейными конфигурациями, которые впоследствии оказываются значимыми для развития научно-технического знания.

\subsection{Абляционный эксперимент}

Цель абляционного эксперимента состоит в оценке вклада отдельных компонентов предложенного метода в итоговое качество ранжирования публикаций по степени их концептуальной новизны. Такая постановка позволяет установить, является ли эффективность метода следствием всей предложенной архитектуры в целом или же она может быть достигнута существенно более простыми средствами.

В качестве полной модели рассматривается исходный вариант метода, включающий: сегментацию текста на смысловые фрагменты, извлечение идейных формулировок сегментов с помощью LLM, построение множества базовых идей, представление публикаций в виде идейных спектров, а также вычисление интегрального показателя концептуальной новизны на основе выбранной комбинации частных метрик, например структурной новизны, редкости идейной конфигурации, когерентности идейного профиля и локальной плотности идейного пространства.

Далее формируется набор упрощённых вариантов метода, в каждом из которых удаляется или заменяется один из ключевых компонентов. В частности, могут быть рассмотрены следующие абляционные конфигурации:
\begin{itemize}
    \item модель без смысловой сегментации текста;
    \item модель без этапа LLM-извлечения идейных формулировок;
    \item модель без кластеризации идей и построения множества базовых идей;
    \item модель, использующая только одну из частных метрик новизны, например только $D_M(w)$ или только $R(w)$;
    \item модель без компонента когерентности $C(w)$;
    \item модель без учёта локальной плотности $\rho(w)$.
\end{itemize}

Пусть
\[
\mathcal{M}_0
\]
обозначает полную модель, а
\[
\mathcal{M}_1,\dots,\mathcal{M}_L
\]
--- её абляционные варианты. Для каждого варианта на одном и том же корпусе и в одной и той же экспериментальной постановке вычисляются значения итогового ранжирующего показателя, после чего оценивается качество выявления прорывных публикаций по тем же метрикам, что и в предыдущем эксперименте.

Если обозначить через
\[
Q(\mathcal{M}_\ell)
\]
обобщённую метрику качества модели $\mathcal{M}_\ell$, то основной интерес представляет сравнение
\[
Q(\mathcal{M}_0) \quad \text{и} \quad Q(\mathcal{M}_\ell), \qquad \ell=1,\dots,L.
\]
Снижение качества при удалении конкретного компонента интерпретируется как свидетельство его значимого вклада в работу всей системы. Напротив, если исключение некоторого компонента практически не изменяет результат, это может указывать на его избыточность либо на необходимость пересмотра способа его использования в модели.

Ожидаемый результат абляционного эксперимента состоит в том, что полная модель $\mathcal{M}_0$ будет демонстрировать наилучшее качество либо, по крайней мере, устойчивое преимущество над большинством упрощённых вариантов. Такой результат позволит эмпирически подтвердить, что эффективность предложенного подхода определяется не отдельной частной эвристикой, а согласованной работой всех его основных этапов: концептуальной сегментации, извлечения идейных структур, построения пространства базовых идей и интегральной оценки новизны публикации.

Таким образом, сочетание ретроспективного эксперимента и абляционного анализа позволяет проверить предложенный метод с двух сторон: с точки зрения его способности выявлять содержательно значимые публикации и с точки зрения внутренней обоснованности его архитектуры.








\bibliography{main}

\end{document}


